\section{Countervaluations and complete-count attainment}
\label{sec:boundaries}

The graph constructions distinguish three tasks: evaluating one supplied
tuple, finding a witness, and certifying that no witness exists. Two
precise consequences express the remaining distinction at a fixed control.

\begin{proposition}[A finite-box characterization]
\label{prop:counter-bound}
Fix the control and parsing policy. Universal weak and strict canonical
comparison on terms with ground free names, of arbitrary result type,
are co-recursively-enumerable. This comparison problem is decidable if
and only if there is a total computable function $b(Q)$ such that
every false query $Q$ has a countervaluation with every raw ground
coordinate at most $b(Q)$.
\end{proposition}
\begin{proof}
Enumerate finite tuples of positive integers and evaluate the supplied
query at each tuple. The closed positive coordinates and all scalar
values are computable, so failure is semidecidable. If the bound is
supplied, search the finite positive box through $b(Q)$ and decide the
query. Conversely, given a decider, return zero on a true query. On a
false query enumerate until a countervaluation is found and return its
largest coordinate. This latter search is entered only when its
termination is known. With no free names there is one empty valuation,
whose largest coordinate is taken to be zero.
\end{proof}

The bound concerns the location of a countertuple, not the time required
to evaluate a given tuple. At the controls of
Theorem~\ref{thm:control-hardness} no such computable bound exists,
even for its restricted fragment. At a control with a proved comparison
decider the proposition supplies a computable bound, without a useful
size estimate.

For the next result specialize to the unchanged standard control $G$.
We give the short residue argument required by the reduction.

\begin{lemma}[Standard-control residues]\label{lem:standard-residues}
For every $m\ge1$ one can compute $r_m\in\{0,\ldots,m-1\}$
such that
\[
 F_j(n)\equiv r_m\pmod m\quad(j\ge1, n\ge m).
\]
Consequently $G(t)\equiv r_m\pmod m$ for
$t\ge\max(0,m-100)$.
\end{lemma}
\begin{proof}
Put $r_1=0$. For $m>1$ write $m=2^a b$ with $b$ odd, and set
$q=\varphi(b)$ if $b>1$, and $q=1$ otherwise. Since $q<m$,
recursively compute $r_q$ and use the Chinese remainder theorem to set
\[
 r_m\equiv0\pmod{2^a},\qquad r_m\equiv2^{r_q}\pmod b.
\]
For the towers $T_0(n)=n$ and $T_{h+1}(n)=2^{T_h(n)}$,
induction on $m$ gives $T_h(n)\equiv r_m\pmod m$ whenever
$h,n\ge m$. The last exponent is at least $n\ge a$, giving
the power-of-two congruence, and has residue $r_q$ modulo $q$ by
the smaller-modulus induction, giving the odd congruence by Euler's
theorem. Apply this to $F_1(n)=T_{n+1}(n)$. At every higher level
all intermediate arguments are at least $n$; the final application
of the preceding level proves the same assertion. The shift by 100
gives the statement for $G$.
\end{proof}

\begin{proposition}[Attainment for a supplied actual unary source]
\label{prop:actual-attainment}
Let $H:o\to o$ be a supplied closed original source at control $G$,
with its actual canonical vector $(c_H,a_H)$. Put
\[
 \lambda=\omega\otimes a_H,\qquad
 C_H=N(\lambda)+c_H+1.
\]
Given $p,m\ge1$ and residues $a,r$, deciding whether there is a
genuine input $t\ge N(\lambda)$ such that
\[
 t\equiv a\pmod p,\qquad D_\lambda(t)\equiv r\pmod m
\]
reduces effectively to actual-source arithmetic satisfiability, and
hence to one universal ground comparison with its answer complemented.
\end{proposition}
\begin{proof}
The literal application identity is
\[
 F_H(n)=\can{HX}_0-1=J_\lambda(C_H+n)-1,
 \qquad x_0=n+1.
\]
Choose an effective $b\ge C_H$ above the threshold in
Lemma~\ref{lem:standard-residues}, with $b\equiv a\pmod p$.
Every terminal state on that tail is a positive iterate of $G$, so
$Z_\lambda(t)\equiv r_m\pmod m$. The complete-count identity
therefore gives
\[
 F_H(t-C_H)\equiv D_\lambda(t)+r_m-1\pmod m.
\]
Let $s$ be the least nonnegative residue of $r+r_m-1$ modulo $m$.
The desired tail query is exactly
\[
 \exists k,y\in\N\quad
       F_H(b-C_H+pk)=my+s.
\]
It is a formula in actual-source arithmetic with one supplied $H$
graph and fixed integer coefficients. Its two independent witnesses
remain shared wherever they occur. The Boolean compiler gives the
stated comparison reduction.

For every genuine $t<b$, compute the exact descent and its complete
event count and test the two congruences. There are finitely many such
inputs. If one succeeds, the whole attainment question has a positive
answer; otherwise the tail sentence decides it. Thus no part of the
genuine finite prefix is omitted.
\end{proof}

The count is complete: it includes every nonfinite event after any
intermediate ordinal cut. The actual scalar $c_H$ and the norm offset
$N(\lambda)$ remain in its observable input ray. This is a necessary
subproblem for a full standard-control comparison decider. It is not
an assertion that every ordinal head has a supplied source, that every
isolated intermediate counter can be observed, or that such attainment
queries are equivalent to effective periodic tables at all auxiliary
heads. The carrier certificates of Section~\ref{sec:head-graph}
enlarge the class of justified graph observations through a separate
construction, with their own stated premises.
